\section*{الفصل \ref{ch:g4:world-physical-map} --- الخريطة الطبيعية للعالم}

\begin{solution}{exo:g4:world-physical-map:1}
\omterm{def:g4:world-physical-map:altitude}{ارتفاع} مكان ما هو علوّه فوق \omterm{def:g4:world-physical-map:altitude}{مستوى سطح البحر}، بالأمتار. ويُقاس من \omterm{def:g4:world-physical-map:altitude}{مستوى
سطح البحر}، أي المستوى المتوسط لسطح \omterm{def:g2:water-on-earth:ocean}{البحر}.
\end{solution}

\begin{solution}{exo:g4:world-physical-map:2}
(أ) أخضر داكن: بين 0 و200 متر. (ب) أصفر فاتح: بين 500 و1\,000
متر. (ج) أبيض تقريبًا: فوق 4\,000 متر.
\end{solution}

\begin{solution}{exo:g4:world-physical-map:3}
(أ) أمريكا الجنوبية. (ب) آسيا. (ج) إفريقيا. (د) أمريكا الشمالية.
\end{solution}

\begin{solution}{exo:g4:world-physical-map:4}
أعلى \omterm{def:g1:landscapes-around-us:landform}{جبل} هو \omterm{def:g1:landscapes-around-us:landform}{جبل} إفرست (8\,849 مترًا). وأكبر \omterm{def:g1:landscapes-around-us:water}{بحيرة} هي \omterm{def:g2:water-on-earth:ocean}{بحر} قزوين؛ وأعمق
\omterm{def:g1:landscapes-around-us:water}{بحيرة} هي \omterm{def:g1:landscapes-around-us:water}{بحيرة} بايكال.
\end{solution}

\begin{solution}{exo:g4:world-physical-map:5}
كلما ادكنّ الأزرق ازداد \omterm{def:g2:water-on-earth:ocean}{البحر} عمقًا: \omterm{def:g2:water-on-earth:ocean}{فالبحر} ضحل قرب \omterm{def:g1:landscapes-around-us:water}{السواحل}
(أقل من 200 متر) وأعمق بكثير في وسط المحيطات (أكثر
من 4\,000 متر في أماكن كثيرة).
\end{solution}

\begin{solution}{exo:g4:world-physical-map:6}
\omterm{def:g4:world-physical-map:physical-map}{الخريطة الطبيعية} تبين \omterm{def:g4:world-physical-map:altitude}{ارتفاع} السطح. وفي القارة القطبية الجنوبية وغرينلاند
السطح هو أعلى طبقة سميكة من الجليد، تعلو كثيرًا فوق \omterm{def:g4:world-physical-map:altitude}{مستوى سطح البحر}،
فتأخذان ألوان الأرض المرتفعة.
\end{solution}

\begin{solution}{exo:g4:world-physical-map:7}
قاع غور تشالنجر (قرابة 11\,000 متر تحت \omterm{def:g4:world-physical-map:altitude}{مستوى سطح البحر})؛ ثم متوسط
قاع \omterm{def:g2:water-on-earth:ocean}{المحيط} (نحو 3\,700 متر تحته)؛ ثم شاطئ (نحو
0 متر)؛ ثم قرية على \omterm{def:g4:world-physical-map:altitude}{ارتفاع} 1\,500 متر؛ ثم قمة \omterm{def:g1:landscapes-around-us:landform}{جبل} إفرست
(8\,849 مترًا).
\end{solution}

\begin{solution}{exo:g4:world-physical-map:8}
يتكون \omterm{def:g1:landscapes-around-us:water}{النهر} من جداول كثيرة تجتمع معًا، والناس لا يتفقون على الجدول الذي
يُعدّ منبعه الحقيقي. فإذا قيس \omterm{def:g1:landscapes-around-us:water}{النهر} من منبع أبعد، كان أطول. والنيل
والأمازون متقاربان في الطول إلى حد كبير (كل منهما بين 6\,700 و6\,800
كيلومتر تقريبًا) حتى إن الجواب يتوقف على ذلك الاختيار.
\end{solution}

\begin{solution}{exo:g4:world-physical-map:9}
ربع 30 مليونًا سبعة ملايين ونصف مليون كيلومتر مربع (30 مقسومًا على 4)؛
والثلث 10 ملايين (30 مقسومًا على 3). والصحراء الكبرى، بمساحة ثمانية ملايين
ونصف مليون، أكثر من ربع إفريقيا لكنها أقل من
ثلثها.
\end{solution}

\begin{solution}{exo:g4:world-physical-map:10}
$2500 - 1000 = 1500$: الكوخ أعلى بمقدار 1\,500 متر. والكوخ، بين
2\,000 و4\,000 متر، بني؛ أما \omterm{def:g1:landscapes-around-us:water}{البحيرة}، على \omterm{def:g4:world-physical-map:altitude}{ارتفاع} 1\,000 متر، فتقع على
الحد بين الأصفر الفاتح (من 500 إلى 1\,000) والبني الفاتح (من 1\,000 إلى
2\,000): وكلا الجوابين مقبول.
\end{solution}

\begin{solution}{exo:g4:world-physical-map:11}
10 كيلومترات تساوي 10\,000 متر. $10000 - 8849 = 1151$: الطائرة أعلى
من \omterm{def:g1:landscapes-around-us:landform}{جبل} إفرست، بمقدار 1\,151 مترًا. وفوق غور تشالنجر، تكون على \omterm{def:g4:world-physical-map:altitude}{ارتفاع}
10\,000 متر فوق \omterm{def:g2:water-on-earth:ocean}{البحر}، وقاع الغور يقع 11\,000 متر تحت \omterm{def:g2:water-on-earth:ocean}{البحر}:
$10000 + 11000 = 21000$ متر فوق قاع
الغور.
\end{solution}
